\documentclass[11pt, a4paper]{article}

\usepackage[a4paper, margin=2.5cm]{geometry}
\usepackage{fontspec}
\usepackage[english, bidi=basic, provide=*]{babel}
\babelprovide[import, onchar=ids fonts]{english}
\babelfont{rm}{Noto Sans}

\usepackage{amsmath, amssymb, amsthm}
\usepackage{booktabs}
\usepackage{tabularx}
\usepackage{xcolor}
\usepackage{enumitem}
\usepackage{listings}
\usepackage{fancyhdr}

\definecolor{accentblue}{RGB}{24, 78, 135}
\definecolor{mutedgray}{RGB}{95, 105, 115}
\definecolor{codebg}{RGB}{246, 248, 250}
\definecolor{rulegray}{RGB}{205, 212, 218}

\usepackage[colorlinks=true, linkcolor=accentblue, citecolor=accentblue, urlcolor=accentblue]{hyperref}

\setlength{\headheight}{14pt}
\pagestyle{fancy}
\fancyhf{}
\fancyhead[L]{\small\color{mutedgray}Comprehensive Sample Document}
\fancyhead[R]{\small\color{mutedgray}\leftmark}
\fancyfoot[C]{\small\color{mutedgray}\thepage}
\renewcommand{\headrulewidth}{0.4pt}

\theoremstyle{plain}
\newtheorem{theorem}{Theorem}[section]
\newtheorem{lemma}[theorem]{Lemma}

\theoremstyle{definition}
\newtheorem{definition}[theorem]{Definition}

\theoremstyle{remark}
\newtheorem{remark}[theorem]{Remark}

\lstset{
  language=Python,
  basicstyle=\ttfamily\small,
  keywordstyle=\color{accentblue}\bfseries,
  commentstyle=\color{mutedgray}\itshape,
  backgroundcolor=\color{codebg},
  frame=single,
  rulecolor=\color{rulegray},
  breaklines=true,
  numbers=left,
  numberstyle=\tiny\color{mutedgray},
  xleftmargin=1.5em,
  framexleftmargin=1.2em,
  aboveskip=1em,
  belowskip=1em
}

\title{\textbf{A Comprehensive Sample Document}\\
\large Demonstrating Structure, Mathematics, Tables, Code, and References}
\author{Jane Doe \and John Smith}
\date{\today}

\begin{document}

\maketitle

\begin{abstract}
This document serves as a comprehensive reference template illustrating the core structural and typographic elements used in academic and technical writing. It covers hierarchical sectioning, inline text formatting, lists, display mathematics, theorem environments, multi-column and wrapping tables, source code listings, internal cross-references, and bibliographic citations.
\end{abstract}

\tableofcontents
\vspace{1em}
\hrule
\vspace{1.5em}

\section{Typography and Document Structure}
\label{sec:typography}

Well-structured technical documents rely on clear hierarchy and consistent typographic conventions. Standard text formatting includes \textbf{boldface} for strong emphasis, \textit{italics} or \emph{emphasized text} for introduced terms, \texttt{monospaced text} for identifiers, and \textsc{Small Capitals} for acronyms or formal titles.

Special punctuation and reserved characters can be typeset cleanly in prose: an en-dash is used for numeric ranges such as pages 12--24, whereas an em-dash---like this one---sets off a parenthetical remark. Reserved symbols such as \&, \%, \$, \#, and \_ are escaped so they appear literally.\footnote{Footnotes provide supplementary commentary without interrupting the primary reading flow.}

\subsection{Block Quotations}
\label{subsec:quotes}

Extended excerpts or notable statements are set apart using the quotation environment:

\begin{quote}
``Premature optimization is the root of all evil (or at least most of it) in programming.'' --- Donald E. Knuth \cite{knuth1984}
\end{quote}

\subsubsection{Paragraph-Level Headings}
\paragraph{Run-in headings} are useful when a section requires fine-grained thematic divisions without cluttering the table of contents.

\section{Lists and Enumerations}
\label{sec:lists}

Structured lists help readers scan procedures, criteria, and definitions quickly.

\subsection{Nested Itemized and Enumerated Lists}

\begin{enumerate}
  \item \textbf{Planning Phase:} Define scope, audience, and structural outline.
  \begin{itemize}
    \item Identify primary data sources and citations.
    \item Determine required tables, equations, and code examples.
  \end{itemize}
  \item \textbf{Drafting Phase:} Write prose with consistent terminology and clear transitions.
  \item \textbf{Review Phase:} Verify all cross-references, mathematical derivations, and bibliographic entries.
\end{enumerate}

\subsection{Description Lists}

\begin{description}
  \item[Precision] The proportion of positive identifications that were actually correct.
  \item[Recall] The proportion of actual positives that were correctly identified.
  \item[$F_1$ Score] The harmonic mean of precision and recall, balancing both metrics into a single value.
\end{description}

\section{Mathematical Typesetting}
\label{sec:math}

Mathematical expressions may appear inline, such as Euler's identity $e^{i\pi} + 1 = 0$, or as numbered display equations for key formulas.

\subsection{Numbered Equations and Aligned Derivations}

The Gaussian probability density function with mean $\mu$ and variance $\sigma^2 > 0$ is given by:
\begin{equation}
  \label{eq:gaussian}
  f(x \mid \mu, \sigma^2) = \frac{1}{\sqrt{2\pi\sigma^2}} \exp\!\left( -\frac{(x - \mu)^2}{2\sigma^2} \right).
\end{equation}

Multi-step derivations are aligned around relational operators using the \texttt{align} environment:
\begin{align}
  \sum_{k=1}^{n} k &= 1 + 2 + \dots + n \nonumber \\
  &= \frac{n(n + 1)}{2}, \label{eq:arith-sum} \\
  \int_{0}^{\infty} e^{-ax}\,dx &= \lim_{b \to \infty} \left[ -\frac{1}{a} e^{-ax} \right]_{0}^{b} = \frac{1}{a} \quad (a > 0). \label{eq:exp-int}
\end{align}

\subsection{Matrices and Piecewise Functions}

Linear transformations can be represented in matrix form. For a rotation by angle $\theta$ in $\mathbb{R}^2$:
\begin{equation}
  \label{eq:rotation}
  R(\theta) =
  \begin{bmatrix}
    \cos\theta & -\sin\theta \\
    \sin\theta & \cos\theta
  \end{bmatrix}.
\end{equation}

Piecewise-defined functions, such as the Rectified Linear Unit (ReLU), use the \texttt{cases} environment:
\begin{equation}
  \label{eq:relu}
  \operatorname{ReLU}(x) =
  \begin{cases}
    x & \text{if } x > 0, \\
    0 & \text{if } x \le 0.
  \end{cases}
\end{equation}

\section{Theorems, Definitions, and Proofs}
\label{sec:theorems}

Formal mathematical writing organizes claims into numbered definitions, lemmas, theorems, and proofs.

\begin{definition}[Convergence of a Sequence]
\label{def:convergence}
A sequence of real numbers $(a_n)_{n=1}^{\infty}$ is said to \emph{converge} to a limit $L \in \mathbb{R}$ if for every $\varepsilon > 0$, there exists an integer $N \ge 1$ such that $|a_n - L| < \varepsilon$ for all $n \ge N$.
\end{definition}

\begin{theorem}[Uniqueness of Limits]
\label{thm:uniqueness}
If a sequence $(a_n)$ converges to $L_1$ and also converges to $L_2$, then $L_1 = L_2$.
\end{theorem}

\begin{proof}
Suppose for contradiction that $L_1 \neq L_2$, and set $\varepsilon = \frac{1}{2}|L_1 - L_2| > 0$. By Definition~\ref{def:convergence}, there exist integers $N_1$ and $N_2$ such that $|a_n - L_1| < \varepsilon$ for all $n \ge N_1$ and $|a_n - L_2| < \varepsilon$ for all $n \ge N_2$. Let $N = \max(N_1, N_2)$. Then for $n \ge N$, the triangle inequality yields
\[
  |L_1 - L_2| \le |L_1 - a_n| + |a_n - L_2| < \varepsilon + \varepsilon = |L_1 - L_2|,
\]
which is a contradiction. Hence $L_1 = L_2$.
\end{proof}

\begin{remark}
Theorem~\ref{thm:uniqueness} holds in any Hausdorff topological space, not only in $\mathbb{R}$ with the standard Euclidean metric.
\end{remark}

\section{Tables and Data Presentation}
\label{sec:tables}

Publication-quality tables avoid vertical rules and rely on horizontal rules of varying thickness (\texttt{booktabs}) to separate headers, body rows, and summary totals.

\subsection{Numeric Comparison Table}

Table~\ref{tab:complexity} summarizes the asymptotic time and space complexity of three classical sorting algorithms.

\begin{table}[htbp]
  \centering
  \caption{Asymptotic time and auxiliary space complexity of selected sorting algorithms.}
  \label{tab:complexity}
  \begin{tabular}{lcccc}
    \toprule
    & \multicolumn{3}{c}{\textbf{Time Complexity}} & \textbf{Auxiliary} \\
    \cmidrule(lr){2-4}
    \textbf{Algorithm} & \textbf{Best} & \textbf{Average} & \textbf{Worst} & \textbf{Space} \\
    \midrule
    Insertion Sort & $O(n)$ & $O(n^2)$ & $O(n^2)$ & $O(1)$ \\
    Merge Sort     & $O(n \log n)$ & $O(n \log n)$ & $O(n \log n)$ & $O(n)$ \\
    Quicksort      & $O(n \log n)$ & $O(n \log n)$ & $O(n^2)$ & $O(\log n)$ \\
    \bottomrule
  \end{tabular}
\end{table}

\subsection{Text-Wrapping Table}

When cells contain multi-line descriptions, \texttt{tabularx} automatically allocates remaining width to \texttt{X} columns, as shown in Table~\ref{tab:environments}.

\begin{table}[htbp]
  \centering
  \caption{Summary of common document environments demonstrated in this template.}
  \label{tab:environments}
  \begin{tabularx}{\textwidth}{lX}
    \toprule
    \textbf{Environment} & \textbf{Primary Purpose} \\
    \midrule
    \texttt{abstract} & Concise summary placed at the opening of an article or report. \\
    \texttt{align} & Multi-line mathematical derivations aligned at specified operators. \\
    \texttt{tabularx} & Full-width tables with automatic line wrapping in designated columns. \\
    \texttt{lstlisting} & Formatted source code blocks with line numbers and syntax styling. \\
    \bottomrule
  \end{tabularx}
\end{table}

\section{Source Code Listings}
\label{sec:code}

Technical reports frequently include short algorithmic implementations. Listing~\ref{lst:quicksort} illustrates a concise recursive implementation of binary search in Python.

\begin{lstlisting}[caption={Iterative binary search returning the index of a target value or $-1$.}, label={lst:quicksort}]
def binary_search(arr, target):
    """Return the index of target in sorted list arr, or -1 if absent."""
    low, high = 0, len(arr) - 1
    while low <= high:
        mid = (low + high) // 2
        if arr[mid] == target:
            return mid
        elif arr[mid] < target:
            low = mid + 1
        else:
            high = mid - 1
    return -1
\end{lstlisting}

\section{Cross-References and Citations}
\label{sec:references}

Internal cross-references remain automatically synchronized as sections, equations, or tables are reordered:
\begin{itemize}
  \item Section~\ref{sec:math} introduces the Gaussian density in Equation~\eqref{eq:gaussian}, the arithmetic series in Equation~\eqref{eq:arith-sum}, and the rotation matrix in Equation~\eqref{eq:rotation}.
  \item Section~\ref{sec:theorems} proves Theorem~\ref{thm:uniqueness} on page~\pageref{thm:uniqueness} using Definition~\ref{def:convergence}.
  \item Section~\ref{sec:tables} compares algorithm performance in Table~\ref{tab:complexity} and lists key environments in Table~\ref{tab:environments}.
\end{itemize}

For comprehensive coverage of document preparation principles and typesetting design, readers are referred to the foundational manuals by Knuth \cite{knuth1984}, Lamport \cite{lamport1994}, and Kopka and Daly \cite{kopka2003}.

\begin{thebibliography}{9}

\bibitem{knuth1984}
Donald E. Knuth.
\textit{The \TeX book}.
Addison-Wesley, Reading, Massachusetts, 1984.

\bibitem{lamport1994}
Leslie Lamport.
\textit{\LaTeX: A Document Preparation System}.
Addison-Wesley, Reading, Massachusetts, 2nd edition, 1994.

\bibitem{kopka2003}
Helmut Kopka and Patrick W. Daly.
\textit{Guide to \LaTeX}.
Addison-Wesley, Boston, 4th edition, 2003.

\end{thebibliography}

\end{document}